\section{Ús de la intel·ligència artificial}

L’equip ha adoptat un conjunt de principis inspirats en les recomanacions de la Unió Europea sobre intel·ligència artificial fiable (\textit{Trustworthy AI}) i en els principis recollits a l’\textit{AI Act} europeu.

Els principis aplicats durant el desenvolupament del projecte són els següents:

\begin{itemize}

\item \textbf{Supervisió humana (Human oversight)}  
Les eines d’intel·ligència artificial s’han utilitzat exclusivament com a assistents tècnics. Totes les decisions d’arquitectura, disseny i implementació han estat preses pels membres de l’equip de desenvolupament.

\item \textbf{Responsabilitat del desenvolupador}  
L’equip assumeix la responsabilitat completa sobre el codi final desenvolupat i sobre les decisions tècniques adoptades en el sistema.

\item \textbf{Verificació del contingut generat}  
Qualsevol suggeriment generat per eines d’intel·ligència artificial ha estat analitzat críticament, validat i, si escau, adaptat pels desenvolupadors abans d’incorporar-se al projecte.

\item \textbf{Transparència en l’ús d’eines}  
L’ús d’eines d’intel·ligència artificial es documenta explícitament en aquesta memòria amb l’objectiu de garantir transparència acadèmica i evitar qualsevol ambigüitat sobre l’autoria del treball realitzat.

\item \textbf{Protecció de dades i privacitat}  
Durant l’ús de models de llenguatge no s’ha introduït informació confidencial, claus criptogràfiques, credencials d’accés ni dades personals reals.

\end{itemize}

\subsection{Eines d’intel·ligència artificial utilitzades}

Durant el desenvolupament del projecte s’han utilitzat diverses eines d’intel·ligència artificial generativa per assistir en tasques d’aprenentatge, documentació tècnica i exploració de solucions.

\subsubsection{Models de llenguatge (LLM)}

S’han utilitzat models de llenguatge de gran mida com a assistents tècnics per a:

\begin{itemize}

\item aclariment de conceptes tècnics complexos
\item exploració de patrons d’arquitectura
\item generació d’esborranys de documentació
\item revisió conceptual de fragments de codi
\item generació d’exemples d’implementació

\end{itemize}

Entre els models utilitzats es troben:

\begin{itemize}

\item \textbf{ChatGPT (OpenAI)}  
Utilitzat com a assistent conversacional per a la comprensió de conceptes tècnics relacionats amb blockchain, smart contracts i desenvolupament d’aplicacions descentralitzades.

\item \textbf{Gemini (Google DeepMind)}  
Utilitzat com a eina complementària per contrastar explicacions tècniques i explorar diferents enfocaments d’implementació.

\end{itemize}

És important destacar que els models de llenguatge s’han utilitzat únicament com a eines de suport conceptual i mai com a substitut del treball de desenvolupament realitzat per l’equip.

\subsubsection{Eines de prototipatge i generació d’interfícies}

Durant les primeres fases del projecte també s’han utilitzat eines de generació de prototips i assistents visuals basats en intel·ligència artificial.

\begin{itemize}

\item \textbf{Antigravity}  
Eina utilitzada per generar prototips ràpids d’interfície i explorar possibles estructures visuals durant la fase de disseny del front-end.

\end{itemize}

Aquestes eines han permès accelerar la fase inicial de disseny sense substituir el desenvolupament manual del sistema final.

\subsection{Limitacions dels models de llenguatge}

Els models de llenguatge presenten limitacions tècniques importants que han de ser considerades durant la seva utilització:

\begin{itemize}

\item poden generar informació incorrecta o incompleta
\item no posseeixen coneixement actualitzat de tots els entorns tecnològics
\item no comprenen completament el context d’un sistema complex
\item no garanteixen la correcció del codi generat

\end{itemize}

Per aquest motiu, qualsevol contingut generat ha estat tractat com un suggeriment inicial i no com una font de veritat tècnica.

\subsection{IA com a eina de suport a l’aprenentatge}

En el context acadèmic d’aquest projecte, les eines d’intel·ligència artificial han contribuït principalment a:

\begin{itemize}

\item accelerar el procés d’aprenentatge de tecnologies emergents
\item facilitar l’exploració d’alternatives arquitectòniques
\item millorar la qualitat de la documentació tècnica
\item donar suport a la resolució de dubtes conceptuals

\end{itemize}

Aquest enfocament reflecteix la realitat actual de l’enginyeria de programari moderna, on els assistents basats en intel·ligència artificial s’estan integrant progressivament en els fluxos de treball dels desenvolupadors.

No obstant això, el valor intel·lectual del projecte, incloent el disseny del sistema, la implementació dels contractes intel·ligents i l’arquitectura final de la solució, correspon íntegrament al treball realitzat pels membres de l’equip de desenvolupament.